\documentclass[12pt]{article}
\usepackage[utf8]{inputenc}
\usepackage[T1]{fontenc}
\usepackage[margin=1in]{geometry}
\usepackage{setspace}
\usepackage{xcolor}
\usepackage[colorlinks=true, urlcolor=blue!60!black, linkcolor=blue!60!black]{hyperref}
\setlength{\parindent}{0em}
\setlength{\parskip}{0.6em}
\renewcommand{\baselinestretch}{1.0}
\begin{document}
\begin{center}
{\LARGE \textit{Craig v. Boren}, 429 U.S. 190 (1976)\par}
\vspace{0.5em}
{\large Case Brief --- Two-Page Edition\par}
\end{center}

\section{Memory Jogger}
Oklahoma's 3.2\% beer law (men 21, women 18) fell, and sex classifications got intermediate scrutiny (\href{https://supreme.justia.com/cases/federal/us/429/190/}{429 U.S. 190}).

\section{Detailed Case Facts}
Oklahoma barred selling 3.2\% beer to males under 21 and females under 18, though it had equalized the age of majority at 18 in 1972. The law barred only sale, not consumption. Curtis Craig, then under 21, and vendor Carolyn Whitener sued; Oklahoma cited drunk-driving statistics and the Twenty-first Amendment. Ruth Bader Ginsburg filed an \textit{amicus} brief.

\section{Procedural History}
A three-judge district court upheld the law; Craig's claim became moot when he turned 21; the Supreme Court reversed.

\section{Judicial Votes}
Reversed, 7--2. Brennan wrote for White, Marshall, Powell, Stevens, and (except Part II-D) Blackmun; Powell, Stevens, and Blackmun concurred; Stewart concurred in the judgment; Burger and Rehnquist dissented.

\section{Holding}
The vendor had third-party standing. Gender classifications must serve important objectives and be substantially related to them; arrest data (2\% of young men versus 0.18\% of young women) were too tenuous a fit. The Twenty-first Amendment does not dilute equal protection.

\section{Analysis of Opinions}
\textbf{Majority:} A 2\% correlation is an ``unduly tenuous `fit,''' and the data were flawed and the law easily evaded. \textbf{Powell:} preferred \textit{Reed}'s ``fair and substantial relation'' test. \textbf{Stevens:} rejected tiers altogether. \textbf{Stewart:} the law was totally irrational. \textbf{Burger:} no vendor standing and only rational-basis review. \textbf{Rehnquist:} the test came ``out of thin air,'' and the relative disparity made the law rational.

\section{Examples: Future Applications}
\textbf{Same-Side:} different motorcycle license ages for men and women, justified by crash statistics, fail. \textbf{Opposite-Side:} a uniform drinking age of 21 for everyone draws only rational-basis review. \textbf{Fence-Sitter:} state-approved higher auto premiums for young men based on strong actuarial data.

\section{Critique}
Rehnquist showed that focusing on the absolute arrest rate proved too much, and originalists say the Clause has no tiered architecture. Feminist scholars like Reva Siegel fault the ``faux precision'' of a middle tier that leaves women with weaker protection than strict scrutiny.

\section{Key Quotations}
\begin{quote}
``To withstand constitutional challenge, previous cases establish that classifications by gender must serve important governmental objectives and must be substantially related to achievement of those objectives.''
\end{quote}
\noindent --- Justice Brennan, Opinion of the Court, 429 U.S. at 197.

\end{document}
